% Part: lambda-calculus
% Chapter: lambda-definability
% Section: lambda-definable-recursive

\documentclass[../../../include/open-logic-section]{subfiles}

\begin{document}

\olfileid{lam}{ldf}{ldr}
\olsection{الدوال \usetoken{S}{lambda definable} عودية}

ليس الحكم أن كل دالة عودية جزئية !!{lambda definable} فحسب؛ بل عكسه
صحيح أيضًا: فكل دالة~!!{lambda definable} عودية جزئية.

\begin{thm}
  \ollabel{thm:lambda-computable} كل دالة جزئية~$f$
  !!{lambda definable} هي دالة عودية جزئية.
\end{thm}

\begin{proof}
  نذكر رسم البرهان دون استيفاء تفاصيله. نبدأ بترقيم حسابي منظم لحدود
  $\lambd$، أي نسند إلى حدود~$\lambd$ أعداد غودل باستعمال ترميز المتتاليات المعتاد
  بقوى الأعداد الأولية. ثم نعرّف دالة عودية جزئية~$\fn{normalize}(t)$:
  تأخذ عدد غودل~$t$ لحد لامبدا، فتردّ عدد غودل لصورته السوية إن وجدت،
  ولا تكون معرّفة إن لم توجد. ونعرّف أيضًا دالتين عوديتين جزئيتين
  $\fn{toChurch}$ و$\fn{fromChurch}$، تنقل الأولى العدد الطبيعي إلى
  عدد غودل لعدد تشيرش الذي يمثله، وتنقل الثانية في الجهة المقابلة.

  وبها نصوغ حسابًا عوديًا جزئيًا لـ~$f$. فهناك حد~$\lambd$~$F$
  !!{lambda define}s~$f$. ولحساب~$f(n_1, \dots, n_k)$ نحسب أولًا الأعداد
  $\fn{toChurch}(n_i)$، ثم نلحقها بـ~$\Gn{F}$، فنحصل على عدد غودل
  للحد~$F \num{n_1}\dots\num{n_k}$. وبعد ذلك نطبق~$\fn{normalize}$
  في هذا العدد. فإن كانت~$f(n_1, \dots, n_k)$ معرّفة، كانت للحد
  $F \num{n_1}\dots\num{n_k}$ صورة سوية، ولا تكون إلا عدد تشيرش؛
  وإن لم تكن معرّفة، لم تكن للحد صورة سوية، فتكون
  \[\fn{normalize}(\Gn{F\num{n_1}\dots\num{n_k}})\]
  غير معرّفة. وأخيرًا نطبق~$\fn{fromChurch}$ في عدد غودل للصورة
  السوية. وهكذا يطابق هذا الحساب العودي الجزئي مواضع تعريف الدالة
  وقيمها معًا.
\end{proof}

\end{document}
